% Generated draft: requires mathematical and source review.
\documentclass[11pt]{article}
\usepackage[margin=25mm]{geometry}
\usepackage{fontspec}
\setmainfont{Latin Modern Roman}
\usepackage{amsmath,amssymb,mathtools,unicode-math}
\setmathfont{Latin Modern Math}
\usepackage{enumitem,needspace,xurl,hyperref}
\setlength{\parindent}{0pt}
\setlength{\parskip}{5pt}
\setlength{\emergencystretch}{3em}
\newcommand{\N}{\mathbb{N}}
\newcommand{\Z}{\mathbb{Z}}
\newcommand{\Q}{\mathbb{Q}}
\newcommand{\R}{\mathbb{R}}
\newcommand{\C}{\mathbb{C}}
\newcommand{\F}{\mathbb{F}}
\newcommand{\A}{\mathbb{A}}
\newcommand{\PP}{\mathbb{P}}
\newcommand{\E}{\mathbb{E}}
\newcommand{\eps}{\varepsilon}
\newcommand{\dd}{\,\mathrm{d}}
\begin{document}
\section*{UnsolvedMath ID: 30005753}
\textbf{Source problem code:} catalogue entry
\subsection{Definitions and mathematical statement}
A Poisson bracket on a commutative algebra is a bilinear, skew-symmetric operation satisfying the Jacobi identity and the Leibniz rule in each argument. A Poisson--Lie group is a group equipped with a Poisson structure for which multiplication is a Poisson map. A Poisson-homogeneous space is a homogeneous space equipped with a Poisson structure such that the action of its Poisson--Lie group is a Poisson map.

For a Lie algebra \(\mathfrak g\), a tensor \(r\in\mathfrak g\otimes\mathfrak g\) satisfies the classical Yang--Baxter equation if
\[
[r_{12},r_{13}]+[r_{12},r_{23}]+[r_{13},r_{23}]=0,
\]
where the subscripts indicate placement in the three tensor factors. In the usual quasitriangular convention, the symmetric tensor \(r+r^{21}\) is invariant under the diagonal adjoint action. The question concerns the quasitriangular solutions covered by the Belavin--Drinfeld classification. The supplied excerpt does not specify the underlying field, Lie algebras, groups, or normalization of these solutions; these parameters must be recovered from the full source before a definitive formulation is possible.

In the usual algebraic interpretation of cluster compatibility, let \(k\) denote the unspecified base field and let \((X,\pi)\) be a Poisson variety with a cluster structure on its rational function field. Compatibility means that, in every extended cluster \((x_1,\ldots,x_m)\), including frozen variables, there are constants \(\omega_{ij}\in k\) such that
\[
\{x_i,x_j\}_{\pi}=\omega_{ij}x_i x_j
\qquad (1\leq i,j\leq m).
\]
This gives the standard log-canonical interpretation of compatibility. The supplied record does not establish whether this is the complete intended requirement, whether generalized cluster structures are allowed, or what regularity and coordinate-ring conditions are imposed.

The proposed assertion is that every Poisson--Lie structure arising from a Belavin--Drinfeld quasitriangular solution admits a compatible cluster structure, and that the assertion extends to the corresponding Poisson-homogeneous brackets in the class intended by the source. The latter class is described in the excerpt only as brackets ``of the kind described above''; the preceding construction is absent. Consequently, this extension cannot be read as a quantification over all Poisson-homogeneous spaces without further evidence.

The source explicitly qualifies the conjecture: the broader compatibility question has a negative answer in general, with an example on \(\mathrm{SL}_2\) attributed to its reference [9]. The affirmative conjecture is restricted to the Belavin--Drinfeld setting, and its extension to the specified Poisson-homogeneous brackets is attributed to reference [10]. The bibliographic identities and precise hypotheses of these references are not supplied. This is therefore a statement-only formulation of the intended conjectural program, with unresolved scope qualifications, rather than a fully specified universal assertion.
\subsection{Short English statement}
The conjecture asks whether Poisson--Lie structures associated with Belavin--Drinfeld quasitriangular solutions, and a particular associated class of Poisson-homogeneous brackets, admit compatible cluster structures. Compatibility ordinarily means that the Poisson brackets of cluster coordinates are log-canonical in every seed. The supplied excerpt leaves the ambient objects, the homogeneous-space class, and additional cluster requirements unspecified. It also records a counterexample to the unrestricted question on \(\mathrm{SL}_2\).
\subsection{Sources}
\begin{itemize}
\item Supplied local UnsolvedMath registry record, permanent ID \(30005753\), title: Belavin--Drinfeld Cluster Correspondence; exact source problem code: \texttt{catalogue entry}. The record labels the problem \texttt{partially\_solved} and the statement \texttt{corrected\_verified}. Its claim that the official report was checked directly is supplied provenance, not verification performed for this draft.
\item Original OWR report: \emph{Cluster Algebras and Its Applications} (2024), Oberwolfach Reports, pp. 69--136, DOI 10.4171/owr/2024/2. \url{https://doi.org/10.4171/owr/2024/2}. The excerpt is attributed to Misha Gekhtman, ``Poisson maps as a tool for constructing cluster structures,'' joint work with Michael Shapiro and Alek Vainshtein.
\item Supplied literature source labeled ``unified exotic-structure paper'': \url{https://arxiv.org/abs/2308.16701}.
\item Supplied literature source labeled ``generalized GLn paper'': \url{https://arxiv.org/abs/2603.12486}.
\item Literature assessment checked at \texttt{2026-08-22}: the supplied curation labels the status partially solved and states that the correspondence has been established for substantial families while the general program remains incomplete. This is supplied dated curation, not an independently verified literature assessment. The background explicitly requests independent verification of time-sensitive status and cited claims.
\item The excerpt cites Belavin and Drinfeld as reference [1], the general counterexample and initial conjecture as reference [9], and the extension to Poisson-homogeneous brackets as reference [10]. Their bibliographic entries are absent from the supplied record; no identities are inferred here.
\end{itemize}
\end{document}
